\begin{helper}[The replacement exit has a free color and is nondefective]
Let $C$ be a \noderef{ScheideReducedConfiguration} for
$(G,k,e_0,\varphi,T)$, let $D$ be its
\noderef{ScheideDefectiveSetup}, let $R$ assign
\noderef{ScheideReplacementData} to each defective color, and let $L$ be
a \noderef{ScheideLastExitPath} certificate for $\alpha$.  Fix a
defective color $\delta$.  Assume explicitly that
$v=L.\mathit{lastInside}$ belongs to
$U_\delta=\{(R\delta).u1,(R\delta).u2\}$ and that the labeled
$(D.\alpha_0,\alpha)$ \noderef{TwoColorIncidenceGraph}s for $\varphi$ and
$\varphi_\delta$ are equal.  If $v\notin V(T)$, then $v$ misses a color
$\beta\in\noderef{TreeFreeColors}(\varphi_\delta,T_\delta)$, where
$T_\delta=(R\delta).\mathit{replacementTree}$, and
\[
 \alpha\notin\noderef{DefectiveColors}
       (G,\varphi_\delta,V(T_\delta)).
\]
\end{helper}
\begin{proof}
Put $r=D.\gammaIndex(\delta)$,
$Z=V(T)\setminus\{C.\mathit{blockVertex}(r,0),
C.\mathit{blockVertex}(r,1)\}$, and
$\gamma=C.\gamma(r)$.  Also put
$X=\noderef{ScheideEnlargedVertexSet}(R)$.  By
\noderef{ScheideReplacementMaximum},
$(T_\delta,e_0,\varphi_\delta)$ is globally maximum.  Applying
\noderef{ScheideGlobalMaximumOrder} shows that $V(T_\delta)$ is
\noderef{ColoredVertexSetElementary} and closed.  The exact vertex-set
field of \noderef{ScheideReplacementData} is
\[
 V(T_\delta)=Z\cup U_\delta.                              \tag{1}
\]
By the definition of $X$ in \noderef{ScheideEnlargedVertexSet}, its first
summand is $V(T)$.  Hence
$Z\subseteq V(T)\subseteq X$.  Its second summand is the union of all
indexed replacement pairs, and $U_\delta$ is the summand indexed by the
present defective color $\delta$; hence $U_\delta\subseteq X$.  Equation
(1) therefore gives
\[
                 V(T_\delta)\subseteq X.                    \tag{2}
\]
In particular, the assumed membership $v\in U_\delta$, together with
(1), puts $v$ in $V(T_\delta)$.

The base and block positions in
\noderef{ScheideStructuralConfiguration}, together with vertex
injectivity in \noderef{TreeSequence}, place $y_0$ in $Z$.  The exit
field of \noderef{ScheideDefectiveSetup}, expanded through
\noderef{ScheideExitVertex}, places $D.v0(\delta)$ in
\noderef{ScheideExitVertices}; \noderef{ScheideStructuralExitsInBase}
then places it in the base and hence in $Z$.  The missing-set formula in
\noderef{ScheideDefectiveSwitchEffect} preserves $D.\alpha_0$ at $y_0$:
away from $D.v0(\delta)$ it fixes the missing set, while at that vertex
it deletes only $\gamma$, and
$D.\alpha_0\ne\gamma$ is a field of
\noderef{ScheideDefectiveSetup}.

Apply \noderef{ScheideGlobalMaximumMissingCounts} to the replacement
triple and $v\in V(T_\delta)$.  Whether or not $v$ is incident with the
uncolored edge, the corresponding alternative gives a color
$\beta\in\noderef{MissingColors}(G,\varphi_\delta,v)$.  Suppose $\beta$
were used on $T_\delta$.  The theorem
\noderef{ScheideReplacementUsedWitness} would give $z\in Z$ also missing
$\beta$.  Here $z\in V(T)$ while $v\notin V(T)$, so $z\ne v$; both lie in
$V(T_\delta)$ by (1).  This contradicts the elementariness of that set.
Consequently $\beta$ is unused.  Since it is missing at a tree vertex,
the definitions of \noderef{TreeMissingColors} and
\noderef{TreeFreeColors} show that it is free on $T_\delta$.

It remains to exclude defectiveness of $\alpha$.  Suppose instead that
$\alpha\in\noderef{DefectiveColors}(G,\varphi_\delta,V(T_\delta))$.
We first put the recorded path into the color order required by an exit
witness.  By \noderef{TwoColorIncidenceGraph}, adjacency in
$I_G(\psi;a,b)$ tests only whether an incident edge has color $a$ or
color $b$.  Commuting this disjunction proves, for every coloring $\psi$,
\[
                 I_G(\psi;a,b)=I_G(\psi;b,a).
\]
Consequently the assumed graph equality gives
\[
 I_G(\varphi_\delta;\alpha,D.\alpha_0)
 =I_G(\varphi_\delta;D.\alpha_0,\alpha)
 =I_G(\varphi;D.\alpha_0,\alpha).                            \tag{3}
\]
Regard the same ordered walk $L.path$ through (3) as a walk
$P_\delta$ in the left-hand graph.  Equality of the graph structures
does not change the vertices, length, vertex at an index, or support of
the walk.  Thus $L.pathSimple$ makes $P_\delta$ simple.  Moreover,
$L.pathSpans$, together with (3) and the component definition
\noderef{TwoColorComponent}, gives
\[
 \{x:x\in\mathit{support}(P_\delta)\}
 =\operatorname{Comp}_{I_G(\varphi_\delta;\alpha,D.\alpha_0)}(y_0).
                                                               \tag{4}
\]

We now verify every field of the required exit witness.  The base-position
argument above gives $y_0\in Z$, so (1) gives
$y_0\in V(T_\delta)$; the missing-set argument above gives
$D.\alpha_0\in\overline{\varphi_\delta}(y_0)$.  The present
contradiction assumption says that $\alpha$ is defective for
$V(T_\delta)$.  The terminal-exclusion field $L.terminalOutside$ of
\noderef{ScheideLastExitPath} says $L.terminal\notin X$, and (2) therefore
gives $L.terminal\notin V(T_\delta)$.  Together with simplicity and (4),
these facts verify all four clauses of
\noderef{ScheideDefectiveChainPath} for
\[
 (T_\delta,\varphi_\delta,\alpha,D.\alpha_0,
   y_0,L.terminal,P_\delta).
\]

Finally take $i=L.lastIndex$.  The field $L.lastInsideAt$ says that the
vertex of $P_\delta$ at $i$ is $v$, and the assumed membership
$v\in U_\delta$, together with (1), says $v\in V(T_\delta)$.  If a later
position contains a vertex-node $w$,
the field $L.afterLastOutside$ says $w\notin X$; by (2), this implies
$w\notin V(T_\delta)$.  Hence all fields in
\noderef{ScheideExitVertex} hold with
\[
 u=y_0,\qquad \gamma=D.\alpha_0,\qquad \delta=\alpha,
 \qquad z=L.terminal,\qquad P=P_\delta,\qquad i=L.lastIndex.
\]
Thus the definition of \noderef{ScheideExitVertices} gives
$v\in\noderef{ScheideExitVertices}(T_\delta,\varphi_\delta)$.
But \noderef{ScheideExitNoFree}, applied to the replacement globally
maximum triple, makes the missing colors at such an exit vertex disjoint
from the free colors.  This contradicts the color $\beta$ found above,
which is both missing at $v$ and free on $T_\delta$.
Hence $\alpha$ is not defective for $T_\delta$.
\end{proof}
